\section{Síntesis y ampliación}

Esta sección condensa todo el video en un marco reutilizable. El informe trimestral sobre modelos grandes del primer trimestre de 2025 puede entenderse como «el regreso de la inteligencia como eje central»: no hay que dejarse desviar por el éxito repentino de un solo producto ni por el lanzamiento de una sola empresa, sino seguir preguntando de dónde proviene el límite superior del modelo, cómo actúa el modelo, cómo la acción genera retroalimentación, cómo la retroalimentación vuelve al aprendizaje, cómo el producto da cabida a la inteligencia y cómo la empresa convierte su capacidad organizacional en estrategia.

\begin{importantbox}{Seis conclusiones centrales}
Primero, el preentrenamiento sigue siendo determinante para el límite superior de la capacidad base y no debe descuidarse por el auge del postentrenamiento. Segundo, Coding es el entorno digital en el que los modelos obtuvieron primero una capacidad de acción verificable. Tercero, Agent es el salto sistémico de la salida de lenguaje a la ejecución de tareas. Cuarto, Online Learning podría convertir el proceso de uso en un nuevo proceso de producción de datos. Quinto, las barreras de entrada de las aplicaciones provienen del entorno, el estado, la retroalimentación y el control de costos. Sexto, la competencia entre China y Estados Unidos se decidirá, en última instancia, por la creación verificable, y no por los relatos que circulan dentro de cada círculo.
\end{importantbox}

\subsection{Tabla de consulta rápida de términos de este episodio}

Esta sección ofrece una consulta rápida para relacionar este episodio con los episodios anteriores y posteriores de la serie de IA e internet de Zhang Xiaojun (张小珺). Hay una continuidad clara entre los productos Agent de EP97 y EP101, la ruta multimodal de EP102, la inteligencia corporizada de EP106/109, el informe técnico sobre Agent de EP110 y los informes trimestrales sobre modelos grandes de EP127/136.

\noindent\begin{tabularx}{\textwidth}{p{0.23\textwidth}p{0.35\textwidth}X}
\toprule
Término/empresa & Significado en este episodio & Cómo darle seguimiento \\
\midrule
Preentrenamiento & Abrir el límite superior intrínseco del base model & Observar si la siguiente generación de modelos muestra capacidades nuevas, y no solo puntuaciones infladas. \\
Reasoning/RL & Reforzar el razonamiento y el desempeño en tareas & Observar si se traduce en la realización de tareas reales. \\
Coding & Entorno de acción digital y «las manos» del modelo & Observar el pago de los desarrolladores, las llamadas a herramientas y la tasa de éxito en tareas de extremo a extremo. \\
Agent & Ciclo cerrado de percepción, herramientas, memoria y ejecución & Observar las tareas largas, la recuperación ante errores y la controlabilidad. \\
Online Learning & Generar retroalimentación durante el uso y aprender de forma continua & Observar si el entorno, la recompensa, la memoria y la auditoría de seguridad han madurado. \\
DeepSeek & Eficiencia de código abierto y señal de los modelos chinos & Observar la capacidad, el costo, el ecosistema y la iteración continua. \\
Manus/Cursor & Contenedores de producto para la capacidad del modelo & Observar las barreras del flujo de trabajo, los ingresos y la economía de los token. \\
\bottomrule
\end{tabularx}

\subsection{Preguntas para el seguimiento}

Esta sección convierte el informe trimestral en una lista de seguimiento. Con estas preguntas, el lector puede juzgar si el conjunto de juicios de Guangmi (广密) se sigue sosteniendo en los trimestres posteriores de 2025.

\begin{enumerate}
\item ¿La siguiente generación de base model mostrará capacidades nuevas evidentes, o dependerá sobre todo del post-training para inflar los resultados en benchmark?
\item En el escenario de Coding, ¿los modelos líderes seguirán siendo los de la serie Anthropic/Claude, o OpenAI, Google, DeepSeek y otros volverán a alcanzarlos?
\item ¿Podrán los productos Agent pasar de la demostración a la entrega estable, en especial en lo que respecta a las tareas largas, los permisos, la seguridad y el costo?
\item ¿Surgirá un paradigma de ingeniería reproducible para Online Learning, en lugar de quedarse en el plano conceptual?
\item Entre las empresas de modelos y las empresas de aplicaciones, ¿quién podrá controlar mejor el estado del usuario, el entorno de la tarea y los datos de retroalimentación?
\item ¿La ruta de eficiencia de código abierto al estilo de DeepSeek podrá seguir extendiéndose a más modelos y ecosistemas de herramientas?
\item En la competencia de IA entre China y Estados Unidos, ¿las restricciones de capacidad de cómputo se compensarán en parte con la eficiencia algorítmica, la optimización de ingeniería y el ecosistema de código abierto?
\end{enumerate}

\subsection{Lecturas adicionales}

\begin{itemize}
\item Si te interesan los productos Agent y el emprendimiento de aplicaciones, puedes contrastar con EP101 YouWare, el informe técnico sobre Agent de EP110 y la historia técnica de Agent de EP139.
\item Si te interesan la multimodalidad y los modelos del mundo, puedes contrastar con la entrevista a Zhang Xiangyu (张祥雨) de EP102 y la entrevista a Xie Saining (谢赛宁) de EP133.
\item Si te interesa la línea de la inteligencia corporizada y la robótica, puedes contrastar con EP106 Wang He (王鹤), EP109 Guanglun Zhineng (光轮智能) y la versión con proyección de pantalla sobre VLA `eiQFomOuCJs`.
\item Para la continuidad de los informes trimestrales sobre modelos grandes, puedes contrastar con EP127, EP136 y las revisiones trimestrales de 2023/2024 que se enumeran en la descripción del video.
\end{itemize}

\end{document}
